\chapter{Surface complexation}
\def\shorttitle{Surface complexation} %

\section{PHT3D Exercise , cape Cod}
The Cape Cod research site is near a wastewater-treatment facility at the Massachusetts Military Reservation (MMR). A treated-sewage plume originated 
from the facility�s infiltration beds, which were used from about 1936 to 1995. The plume extends more than 6 kilometers from the disposal site in the 
sand and gravel aquifer and contains a complex mixture of phosphate, nitrate, metal ions (such as zinc), detergents, organic chemicals, and microbes. 
The plume and adjacent parts of the aquifer  
serve as an ideal field laboratory where a multidisciplinary team of scientists is investigating how contaminants are transported in groundwater.

\subsection{Groundwater Flow}
Regionally,  groundwater  flow  is  to  the  south,  but  locally  the 
flow  direction  is  from  the  sewage  treatment  facility  toward  the 
eastern  shore  of  Ashumet  Pond  (Figure  1).  Groundwater  flow 
is  nearly  horizontal.  Approximately  25  cm/yr  of areal  re- 
charge  results  in a small  vertical  component  to  flow  as  well as 
in  the  accumulation  of a  lens  of  pristine  groundwater  on  top  of 
the sewage  plume,  the thickness  of which  increases  with dis- 
tance  away  from the sewage  treatment  plant [LeBlanc,  1984; 
LeBlanc  et  al., 1991].  The magnitude  of  the  hydraulic  gradient 
is  approximately  1.7  m/km; seasonal  fluctuations  in  the  flow 
system  result  in  small  variations  (approximately  0.2  m /km)  in
the magnitude  of the  gradient  as  well as  in  a 15  �  annual  fluc- 
tuation  in flow  direction.

\subsection{Chemistry}
Contrasting  chemical  conditions  between  sewage-contami- 
nated  and  ambient  groundwater  gave  rise  to longitudinal  and 
vertical  gradients  in  groundwater  chemistry.  The  ambient 
groundwater  surrounding  the  sewage  plume  had  high  concen- 
trations  of dissolved  oxygen,  pH  values  in  the range  5.3-5.8, 
and  low  concentrations  of dissolved  salts  [LeBlanc  et  al., 1991; 
Hess  et al.,  1996;  Kent and Maeder, 1999].  Sewage-contami- 
nated  groundwater  under  the disposal  beds  had  high  concen- 
trations of  nitrate;  concentrations  of  dissolved  oxygen  de- 
creased  with depth [Smith  et al.,  1991a].  Biodegradation  of 
organic  matter  in  conjunction  with  solute  transport  away  from 
the  disposal  beds  resulted  in  the  depletion  of dissolved  oxygen 
throughout  the  plume  and  the  depletion  of  nitrate,  followed  by 
the  accumulation  of  dissolved  Fe(II),  in  the  core  of  the  plume 
[Smith  et  al., 1991b;  Kent  et  al.,  1994].  These  biodegradation 
reactions  influenced  pH  values  in the  sewage  plume 
[Abrams  et  al., 1998;  Abrams  and  Loague,  2000b].

The  principal  chemical  process  influencing  Zn transport  is  assumed  to  be  adsorption  onto 
the  sediments.  A single  reaction  with  the  stoichiometry  adequately  describes  the  pH  effect  observed  in labora- 
tory  batch  studies  of Zn adsorption  onto  the  sediments  [Davis 
et  al., 1998].  Adding  a second  surface  site  with the  same  reaction  stoichiometry  results  in  an  SCM that  fits  both  the  pH and 
Zn-concentration  effect  [Davis  et  al., 1998]. 
Therefore a two-site approach will be applied here
 In  the  two-site  SCM a small  percentage  of the  total 
adsorption  sites  are  considered  to  have  a  much  greater  affinity 
for  adsorption  ("strong  sites,"  >SsOH)  than  the rest  of  the 
sites  ("weak  sites,"  >SwOH): 
Zn  2+  q-  >SsOH  =  >SsOZn  +  +  H  +,  log  SQssozn   
Zn  2+  +  >SwOH-  >SwOZn  +  + H  +,  log  SQswozn.  

The sorption is depedent on pH and the pH varies vertically mainly linked to the intrusion of organic matter
within the cross section. Therfore a typical kinetic model of organic matter degradation is introduced :

....

this kinetic degradation will consume progressively the existing electron acceptors and lead to a reduced pH
in the injection zone. In turn this lower pH may influence the Zinc soprtion. The overall process is tested 
along this case

\subsection{Solute Transport}

The 

\begin{table}[hb]
\caption{Flow and solute transport parameters used for PHT3D Exercise 7}
\begin{center}
\begin{tabular}{lc} 
\toprule
Parameter                                            &         Value \\
\midrule
Flow simulation type                                 &        steady state \\
Simulation time (days)                               &  21000           \\
Time step (days)                               &  60           \\
Model \  extent\ column\ direction\ $ (m)$           &  400            \\
Model \ top\ elevation\ $(m)$                        &   16            \\
Model \ bottom\ elevation\ $(m)$                        &   0            \\
Horizontal hydraulic conductivity\ $(m/day)$         &   110     \\ 
Porosity, $\theta$                                   &           0.39  \\
Piezometric head downstream boundary\ $(m) $         &  14            \\ 
Longitudinal dispersivity, $\alpha_L$   $(m)$             &           1.0 \\
Horizontal transverse dispersivity, $\alpha_{TH}$    &           0.1 \\
Vertical transverse dispersivity,   $\alpha_{TV}$    &           0.01 \\
Effective diffusion coefficient, D$^{\star}$          &             0  \\
\bottomrule

\end{tabular}
\label{single_spec_table_1}
\end{center}
\end{table}

\subsection{Setting up the flow model}

To set up the flow model proceed as follows:

\begin{itemize}

\item In iPHT3D, create a new model and save it in a new and separate directory (folder).

\item Select Parameters $\rightarrow$ 1.Model $\rightarrow$ Model to specify the model dimension ('Xsection')
and type ('free')
\item Select Parameters $\rightarrow$ 1.Model $\rightarrow$ Domain to specify
the domain as described in Table \ref{single_spec_table_1}. For the first simulation select a spacing of 20 m in x direction
and 1 m in y direction

\item Select Parameters $\rightarrow$ Time $\rightarrow$ and set the total Simulation Time to 21000 days. Then define the step size to 60 days

\item Set the boundary conditions for heads by selecting Spatial Attributes $\rightarrow$ Modflow $\rightarrow$ bas.3 boundary conditions corresponding to
IBOUND (Modflow). Use the zone tool to create a zone at the outflow from 0 to 13.9 m

\item to fix heads at the boundaries, use Spatial Attributes $\rightarrow$ Modflow $\rightarrow$ bas.5 Initial heads. provide a value of 16 in the dialog box.
this value will be used everywhere as a starting head and will be used as a BC at places where the BC where fixed (-1) previously. In order to 
set the outflow boundary, draw a line at the same place as the outflow BC and attribute a value of 14 m.

\item the inlet is done by a well between 0 and 15 m height with an total injection rate of 2.16 $m^3/d$

\item select Parameters $\rightarrow$ 2.Flow $\rightarrow$ Parameters and in this dialog set several parametters. 
In lpf.2 select convertible for the type of layer. Use lpf.8 to set the hydraulic conductivity.

\item to fix recharge, use Spatial Attributes $\rightarrow$ Modflow $\rightarrow$ rch.2 recharge. create a zone 
on the top cell with a value of 0.0014 in the dialog box.

\item for porosity value, Select Spatial Attributes $\rightarrow$ MT3DMS $\rightarrow$ btn.11 and set the value to 0.39. 

\end{itemize}
save the model changes.

\subsection{Running MODFLOW}

You have now completed to set up the flow model and you are ready to run MODFLOW.
This step will provide MT3DMS with the flow field that is used to compute 
advective-dispersive transport.

\begin{itemize}

\item Select Parameters $\rightarrow$ 2.Flow $\rightarrow$ Write

\item Select Parameters $\rightarrow$ 2.Flow $\rightarrow$ Run...

\item When the calculation has finished, a window appears, verify that 'NORMAL TERMINATION OF MODFLOW' is written.

\end{itemize}

\subsection{Visualising hydraulic heads}

\begin{itemize}

\item Visualise results by selecting Results $\rightarrow$ 2.Flow $\rightarrow$ Head. The time step can be changed
in Results $\rightarrow$ 1.Model $\rightarrow$ Time.

\end{itemize}

\subsection{Setting up the transport model}

to be copied from previous models

\begin{itemize}

\item Define the observation well Spatial Attributes $\rightarrow$ OBSERVATION $\rightarrow$ obs.1 observations
use the zone tool to draw a point at any location

\item visualize the head contours in Results $\rightarrow$ 3.Transport $\rightarrow$ Tracer. You can change the Time above 
to see another time. It is possible, after selecting one time to use the wheel mouse to stream along time.

\item Visualise the breakthrough curve by selecting Results $\rightarrow$ 5.Observation $\rightarrow$ Breakthrough, then below, 
select the name of the point you just create. A dialog appears that allows you to select value and a breakthrough curves appears.
The curve can be exported to be used in a spreadsheet.

\end{itemize}

\subsection{Chemistry}

in this section we will define the initial composition and solutions used as source with the following steps :
\item import the database
\item the phase in presence everywhere (background) is goethite at 0.1 mol/L. FeS(ppt) must also be ticked but with a null 
concentration (only precipitaiton allowed)
\item As previously the initial solution composition are entered in the Solution Tab, we need the background and two solutions.
the species needed to be activated are Cl,C(4), Fe(2) and Fe(3), Na, O(0), S(-2) and S(6), Orgc, Zn and like always pH and pe. 
Concentrations are given in table ??
\item the solutions are entered in the ph.4 source/sink line at inflow cells as zones as follows :
 from 14 to 16 m : solution 2, from 3 to 14 m : solution 1
from 0 to 3 m : background (solution 0).
\item the recharge composition is entered in ph.5 line at the top cells using solution 3.
\item in Rates specify the degradation rate parameters of Orgc (0,1e-11, 1e-12, 1e-12,1e-17,1e-17 respectively). 
The formula is Orgc -1.0 CH2O(NH3)0.15 1.0.
\item use two surface sites : S_s and S_w with 1.59e-5 and .0018 respective background concentrations

\begin{table}[hb]
\caption{Chemical composition for PHT3D Exercise 7}
\begin{center}
\begin{tabular}{lc} 
\toprule
Parameter  & Background & Solution 1 & Solution 2 & Solution 3\\
\midrule
C(4)      &   1e-4 &   1e-4  &   1e-4 &   1e-4\\
Cl      &   0 &   0  &  0  &   1e-4 \\
Fe(2)      &   0 &   0  &  0  &   0 \\
Fe(3)      &   0 &   0  &  0  &   0 \\
Na      &   2.3e-4 &   2.3e-4 &   2.3e-4 &   3.13e-4 \\
O(0)      &   5e-4 &   0 &   5e-4 &   5e-4\\
Orgc      &   0 &   1e-4 &   0  &  0 \\
S(6)      &   1e-4 &   1e-4  &   1e-4 &   1e-4\\
S(-2)      &   0 &   0  &  0  &   0 \\
Zn      &   0 &   8e-6  &  8e-6  &   0 \\
pH      &   6 &   6 &  6  &   5.6 \\
pe      &   14.6 &   4 &  14.6  &   14.2 \\
\bottomrule

\end{tabular}
\label{spec_table}
\end{center}
\end{table}

